\newcommand{\R}{\mathbb{R}}
\newcommand{\N}{\mathbb{N}}
\newcommand{\Z}{\mathbb{Z}}
\newcommand{\E}{\mathbb{E}}
\newcommand{\Symmetric}{\mathbb{S}}
\newcommand{\diff}{\mathrm{d}}
\newcommand{\D}{\mathrm{d}}
\newcommand{\cov}[1]{\mathrm{Cov}[#1 ]}
\let\P\relax
\newcommand{\P}[1]{\mathbb{P}\left[#1\right]}
\let\E\relax
\newcommand{\E}[1]{\mathbb{E}\left[#1\right]}
\newcommand{\textr}{\textcolor{red}}
\newcommand{\textb}{\textcolor{blue}}
\newcommand{\suchthat}{\mathrm{s.t.}\quad}
\newcommand{\normal}{\mathcal{N}}
\newcommand{\indicator}{\mathds{1}}
\newcommand{\chol}[1]{(#1)^{\frac{1}{2}}}
\newcommand{\maxeig}[1]{\lambda_{\max} \left\{#1 \right\}}
\newcommand{\blkdiag}[1]{\mathrm{blkdiag}(#1)}
\newcommand{\diag}[1]{\mathrm{diag}(#1)}
\DeclareMathOperator*{\minimize}{minimize}


\newtheorem{theorem}{Theorem}
\newtheorem{assumption}{Assumption}
\newtheorem{problem}{Problem}
\newtheorem{remark}{Remark}

\usepackage{cleveref}
\crefname{align}{Eq.}{Eqs.}
\crefname{equation}{Eq.}{Eqs.}
\crefname{figure}{Fig.}{Figs.}
\crefname{table}{Table}{Tables}
\crefname{theorem}{Theorem}{Theorems}
\crefname{definition}{Definition}{Definitions}
\crefname{lemma}{Lemma}{Lemmas}
\crefname{remark}{Remark}{Remarks}
\crefname{assumption}{Assumption}{Assumptions}
\crefname{proof}{Proof}{Proofs}
\crefname{algorithm}{Algorithm}{Algorithms}
\crefname{problem}{Problem}{Problems}
\crefname{proposition}{Proposition}{Propositions}
\crefname{corollary}{Corollary}{Corollaries}
\crefname{section}{Section}{Sections}


\allowdisplaybreaks[1] %allows equations/align to span multiple pages

\newcommand\blfootnote[1]{%
  \begingroup
  \renewcommand\thefootnote{}\footnote{#1}%
  \addtocounter{footnote}{-1}%
  \endgroup
}

\def\showChanges{0} % Change this value to 1 (red) or 0 (black) to switch highlighting condition

% Define the highlighting command
\newcommand{\highlight}[1]{%
    \ifnum\showChanges=1
        \textcolor{red}{#1}%
    \else
        #1%
    \fi
}
% highlighting command for math mode, to preserve spacing
\newcommand{\mathhl}[1]{
    \ifnum\showChanges=1
        \mathcolor{red}{#1}
    \else
        #1
    \fi
}
\newcommand{\needsfix}[1]{%
    \textcolor{blue}{#1}%
}

\usepackage{marginnote}
\usepackage{tcolorbox}
\setlength{\marginparwidth}{0.8in}
\newcommand{\margincomment}[1]{%
    \ifnum\showChanges=1
    \textcolor{blue}{\setlength{\baselineskip}{10pt}\footnotesize\raggedright\marginnote{#1}}%
    \else
    \fi
}